
\section{$[0,1]\times C$ 模型的唯一性}
\label{sec:4-4}

本节证明定理 4.2.34；该定理在研究滤过 $A_\infty$ 代数的同伦等价时
起着基础作用。

\subsection{数目滤过归纳 I}
\label{subsec:4-4-1}
\liftIndex{I094}

在第 4.4.1--4.4.4 小节中，先证明无滤过 $A_\infty$ 代数的情形。
具体考虑如下情形。令 $\overline C_1,\overline C_2$ 为
$A_\infty$ 代数，令 $\overline{\mathfrak C}_1$、
$\overline{\mathfrak C}_2$ 分别为
$[0,1]\times\overline C_1$、$[0,1]\times\overline C_2$ 的模型。
再令
$\overline{\mathfrak f}\colon\overline C_1\to\overline C_2$
为 $A_\infty$ 同态。我们将构造
$A_\infty$ 同态
$\overline{\mathfrak F}\colon
\overline{\mathfrak C}_1\to\overline{\mathfrak C}_2$，使得
\begin{align}
 \widehat{\overline{\mathrm{Eval}}}_{s=s_0}
 \circ\widehat{\overline{\mathfrak F}}
 &=\widehat{\overline{\mathfrak f}}
 \circ\widehat{\overline{\mathrm{Eval}}}_{s=s_0},
 &&s_0=0,1,
 \tag{4.4.1.1}\label{eq:4-4-1-1}\\
 \widehat{\overline{\mathrm{Incl}}}
 \circ\widehat{\overline{\mathfrak f}}
 &=\widehat{\overline{\mathfrak F}}
 \circ\widehat{\overline{\mathrm{Incl}}}.
 \tag{4.4.1.2}\label{eq:4-4-1-2}
\end{align}
这里，帽号表示 bar 复形上诱导的余代数同态，见 (3.2.6)。若
\eqref{eq:4-4-1-1}、\eqref{eq:4-4-1-2} 成立，就称
$\overline{\mathfrak F}$ 在 $\overline{\mathfrak f}$ 之上与
$\overline{\mathrm{Incl}}$、$\overline{\mathrm{Eval}}$ 相容。

回忆，$A_\infty$ 同态 $\overline{\mathfrak F}$ 由一列映射
\[
 \overline{\mathfrak F}_i\colon
 B_i\overline{\mathfrak C}_1[1]
 \longrightarrow\overline{\mathfrak C}_2[1],
 \qquad i=1,2,3,\ldots,
\]
组成，并满足定义 3.2.7 的条件。若
$\overline{\mathfrak F}_i$（$i=1,\ldots,k$）已经定义，而且
\eqref{eq:4-4-1-1}、\eqref{eq:4-4-1-2} 在
$\bigoplus_{i=1}^kB_i\overline{\mathfrak C}_1[1]$ 上成立，也称它与
$\overline{\mathrm{Incl}}$、$\overline{\mathrm{Eval}}$ 相容。

$\overline{\mathfrak F}$ 将通过对数目滤过作归纳来构造。也就是说，按
$k$ 归纳构造
\[
 \overline{\mathfrak F}_k\colon
 B_k\overline{\mathfrak C}_1[1]
 \longrightarrow\overline{\mathfrak C}_2[1].
\]
归纳从下面的结论开始。

\begin{lemma}{4.4.2}\label{lem:4-4-2}
存在与 $\overline{\mathrm{Incl}}$、$\overline{\mathrm{Eval}}$ 相容的
上链同态
\[
 \overline{\mathfrak F}_1\colon
 \overline{\mathfrak C}_1\longrightarrow\overline{\mathfrak C}_2.
\]
\end{lemma}

\begin{proof}
先注意，存在满足 \eqref{eq:4-4-1-2}（在
$B_1\overline C_1=\overline C_1$ 上）的上链映射
$\overline{\mathfrak F}'_1$，但它未必满足
\eqref{eq:4-4-1-1}。例如可以置
\[
 \overline{\mathfrak F}'_1
 =\overline{\mathrm{Incl}}_1
 \circ\overline{\mathfrak f}_1
 \circ(\overline{\mathrm{Eval}}_{s=1})_1.
\]
现在置
\[
 \overline{\operatorname{Err1}}_{s_0}
 =(\overline{\mathrm{Eval}}_{s=s_0})_1
 \circ\overline{\mathfrak F}'_1
 -\overline{\mathfrak f}_1
 \circ(\overline{\mathrm{Eval}}_{s=s_0})_1
\]
以及
\[
 \overline{\operatorname{Err1}}
 =\bigl(\overline{\operatorname{Err1}}_0,
 \overline{\operatorname{Err1}}_1\bigr)
 \in
 \operatorname{Hom}(\overline{\mathfrak C}_1,\overline C_2)
 \oplus
 \operatorname{Hom}(\overline{\mathfrak C}_1,\overline C_2).
\]
于是 $\overline{\operatorname{Err1}}$ 是上链同态。换言之，
$\delta_1\overline{\operatorname{Err1}}=0$；这里 $\delta_1$ 是
$\operatorname{Hom}(\overline{\mathfrak C}_1,\overline C_2)
\oplus\operatorname{Hom}(\overline{\mathfrak C}_1,\overline C_2)$
上的上边界算子，由 $\overline{\mathfrak C}_1$、$\overline C_2$ 上的
$\widehat{\overline{\mathfrak m}}_1$、
$\overline{\mathfrak m}_1$ 诱导。对 $s_0=0,1$，有
\[
 \begin{aligned}
 \overline{\operatorname{Err1}}_{s_0}
 \circ\overline{\mathrm{Incl}}
 &=(\overline{\mathrm{Eval}}_{s=s_0})_1
 \circ\overline{\mathfrak F}'_1
 \circ\overline{\mathrm{Incl}}\\
 &\quad{}-\overline{\mathfrak f}_1
 \circ(\overline{\mathrm{Eval}}_{s=s_0})_1
 \circ\overline{\mathrm{Incl}}
 =0,
 \end{aligned}
\]
因为 $\overline{\mathfrak F}'_1$ 满足
\eqref{eq:4-4-1-2}。下面这个引理将在本节反复使用。

\begin{lemma}{4.4.3}\label{lem:4-4-3}
令 $(\overline D_j,d)$（$j=1,2,3$）为 $R$ 上的上链复形，令
$i\colon\overline D_1\to\overline D_2$ 为上链同态。假设 $i$ 是
上链同伦等价，并且作为 $R$-模同态是分裂单射。

若 $A\in\operatorname{Hom}_R(\overline D_2,\overline D_3)$ 满足
$dA=0$、$A\circ i=0$，则存在
$B\in\operatorname{Hom}_R(\overline D_2,\overline D_3)$，使得
$dB=A$ 且 $B\circ i=0$。
\end{lemma}

\begin{proof}
由假设，
\[
 i^*\colon
 H^*\bigl(\operatorname{Hom}_R(\overline D_2,\overline D_3)\bigr)
 \longrightarrow
 H^*\bigl(\operatorname{Hom}_R(\overline D_1,\overline D_3)\bigr)
\]
是同构。因此存在
$B'\in\operatorname{Hom}_R(\overline D_2,\overline D_3)$，使得
$dB'=A$。由假设，$d(B'\circ i)=0$。所以存在
\[
 \alpha\in\operatorname{Hom}_R(\overline D_2,\overline D_3),
 \qquad
 \beta\in\operatorname{Hom}_R(\overline D_1,\overline D_3),
\]
使得
\[
 B'\circ i=\alpha\circ i+d\beta,
 \qquad d\alpha=0.
\]
由于 $i$ 是分裂单射，存在
$\widetilde\beta\in
\operatorname{Hom}_R(\overline D_2,\overline D_3)$，使
$\widetilde\beta\circ i=\beta$。容易看出
\[
 B=B'-\alpha-d\widetilde\beta
\]
具有所需性质。
\end{proof}

回到引理 4.4.2 的证明。把引理 4.4.3 用于
$A=\overline{\operatorname{Err1}}$，得到
$\overline{\operatorname{Corr}}'$，使得
\[
 \overline{\operatorname{Err1}}
 =\delta\overline{\operatorname{Corr}}',
 \qquad
 \overline{\operatorname{Corr}}'
 \circ\overline{\mathrm{Incl}}_1=0.
\]
由 (4.2.2.4) 和注 4.2.4，同态
\[
 (\overline{\mathrm{Eval}}_{s=0})_1
 \oplus(\overline{\mathrm{Eval}}_{s=1})_1
 \colon\overline{\mathfrak C}_2
 \longrightarrow\overline C_2\oplus\overline C_2
\]
是分裂满射；而
$\overline{\mathrm{Incl}}_1\colon
\overline C_1\to\overline{\mathfrak C}_1$ 是分裂单射。因此存在
$\overline{\operatorname{Corr}}\in
\operatorname{Hom}(\overline{\mathfrak C}_1,
\overline{\mathfrak C}_2)$，使得
\begin{align}
 \bigl((\overline{\mathrm{Eval}}_{s=0})_1,
 (\overline{\mathrm{Eval}}_{s=1})_1\bigr)
 \circ\overline{\operatorname{Corr}}
 &=\overline{\operatorname{Corr}}',
 \tag{4.4.4.1}\label{eq:4-4-4-1}\\
 \overline{\operatorname{Corr}}
 \circ\overline{\mathrm{Incl}}_1&=0.
 \tag{4.4.4.2}\label{eq:4-4-4-2}
\end{align}
现在置
\[
 \overline{\mathfrak F}_1
 =\overline{\mathfrak F}'_1
 -\delta\overline{\operatorname{Corr}}.
\]
容易看出，它具有所需性质。
\end{proof}

至此完成了归纳的第一步。为继续归纳，下面引入 $A_K$ 结构。

\subsection{$A_K$ 结构与同态}
\label{subsec:4-4-2}

定义 $A_K$ 结构还需要一些记号。令 $\overline C$ 为分次 $R$-模，令
$m_1<m_2$ 为自然数。置
\begin{equation}
 B_{m_1\cdots m_2}\overline C[1]
 \cong
 \frac{\displaystyle\bigoplus_{k=1}^{m_2}B_k\overline C[1]}
 {\displaystyle\bigoplus_{k=1}^{m_1-1}B_k\overline C[1]}.
\tag{4.4.5}\label{eq:4-4-5}
\end{equation}
考虑次数为 $1$ 的一列同态
\[
 \overline{\mathfrak m}_k\colon
 B_k\overline C[1]\longrightarrow\overline C[1],
 \qquad k=1,\ldots,K.
\]
由 (3.2.1)，它们诱导余导子
\[
 \widehat{\overline{\mathfrak m}}_k\colon
 B_{1\cdots K}\overline C[1]
 \longrightarrow B_{1\cdots K}\overline C[1].
\]
置
\[
 \widehat{\overline d}_{1\cdots K}
 =\sum_{k=1}^K\widehat{\overline{\mathfrak m}}_k.
\]

\begin{definition}{4.4.6}\label{def:4-4-6}
若
\[
 \widehat{\overline d}_{1\cdots K}
 \circ\widehat{\overline d}_{1\cdots K}=0,
\]
就称 $\overline{\mathfrak m}
=\{\overline{\mathfrak m}_1,\ldots,\overline{\mathfrak m}_K\}$
在 $\overline C$ 上定义了一个 $R$ 上的 $A_K$ 代数结构
\liftIndex{I158}。
\end{definition}

\begin{remark}{4.4.7}\label{rem:4-4-7}
本节不试图把 $A_K$ 代数的概念推广到滤过 $A_\infty$ 代数。这是因为
$\mathfrak m_0$ 不保持 $B_{1\cdots K}C[1]$，相应定义更为复杂。第
7.2.6 小节将在实际需要之处引入滤过 $A_{n,K}$ 代数。
\end{remark}

令 $(\overline C_i,\overline{\mathfrak m})$（$i=1,2$）为
$A_K$ 代数。考虑次数为 $0$ 的一列 $R$-模同态
\[
 \overline{\mathfrak f}_k\colon
 B_k\overline C_1[1]\longrightarrow\overline C_2[1],
 \qquad k=1,2,3,\ldots,K.
\]
按 (3.2.6)，它们诱导映射
\[
 \widehat{\overline{\mathfrak f}}_k\colon
 B_{1\cdots K}\overline C_1[1]
 \longrightarrow B_{1\cdots K}\overline C_2[1].
\]
置
\[
 \widehat{\overline{\mathfrak f}}_{1\cdots K}
 =\sum_{k=1}^K\widehat{\overline{\mathfrak f}}_k.
\]

\begin{definition}{4.4.8}\label{def:4-4-8}
若
\[
 \widehat{\overline{\mathfrak f}}_{1\cdots K}
 \circ\widehat{\overline d}_{1\cdots K}
 =\widehat{\overline d}_{1\cdots K}
 \circ\widehat{\overline{\mathfrak f}}_{1\cdots K},
\]
就称
$\overline{\mathfrak f}
=\{\overline{\mathfrak f}_1,\ldots,\overline{\mathfrak f}_K\}$
是从 $\overline C_1$ 到 $\overline C_2$ 的 $A_K$ 同态
\liftIndex{I159}。\translatornote{原书把每个
$\widehat{\overline{\mathfrak f}}_k$ 都称为“余代数同态”；但按
(3.2.6)，由各分量合成的截断映射才满足余代数同态关系，单个分量在这里
只应称为映射。}两个 $A_K$ 同态的复合用与定义 3.2.8 相同的方式定义。
$A_K$ 同态之间的 $A_K$ 同伦\liftIndex{I160} 以及 $A_K$ 同伦等价
\liftIndex{I161}，也用与第 4.2.2 小节类似的方式定义。
\end{definition}

注意，对 $K'\leq K\leq\infty$，一个 $A_K$ 结构会诱导一个
$A_{K'}$ 结构。下面两个引理将在下一小节使用。

\begin{lemma}{4.4.9}\label{lem:4-4-9}
若 $(\overline C,\overline{\mathfrak m})$ 是 $A_K$ 代数，且
$m_2-m_1<K$，则 $\overline{\mathfrak m}$ 诱导同态
\[
 \widehat{\overline d}_{m_1\cdots m_2}\colon
 B_{m_1\cdots m_2}\overline C[1]
 \longrightarrow B_{m_1\cdots m_2}\overline C[1],
\]
满足
\[
 \widehat{\overline d}_{m_1\cdots m_2}
 \circ\widehat{\overline d}_{m_1\cdots m_2}=0.
\]
\end{lemma}

\begin{lemma}{4.4.10}\label{lem:4-4-10}
若
$\overline{\mathfrak f}\colon\overline C_1\to\overline C_2$
是 $A_K$ 同态
\translatornote{原书在本引理首句的
$\overline{\mathfrak f}$ 上多排了帽号；定义 4.4.8 把基本
$A_K$ 同态记为 $\overline{\mathfrak f}$，帽号只用于它在截断 bar 模上的
诱导映射。}，且 $m_2-m_1<K$，则它诱导同态
\[
 \widehat{\overline{\mathfrak f}}_{m_1\cdots m_2}\colon
 B_{m_1\cdots m_2}\overline C_1[1]
 \longrightarrow B_{m_1\cdots m_2}\overline C_2[1],
\]
满足
\[
 \widehat{\overline{\mathfrak f}}_{m_1\cdots m_2}
 \circ\widehat{\overline d}_{m_1\cdots m_2}
 =\widehat{\overline d}_{m_1\cdots m_2}
 \circ\widehat{\overline{\mathfrak f}}_{m_1\cdots m_2}.
\]
\end{lemma}

证明都很容易，故略去。

\subsection{数目滤过归纳 II}
\label{subsec:4-4-3}
\liftIndex{I094}

现在对 $K$ 作归纳来证明下面的命题。考虑定理 4.2.34 的情形。

\begin{proposition}{4.4.11}\label{prop:4-4-11}
对任意给定的、在 $\overline{\mathfrak f}$ 之上与
$\overline{\mathrm{Incl}}$、$\overline{\mathrm{Eval}}$ 相容的
$A_{K-1}$ 同态 $\overline{\mathfrak F}^{(K-1)}$，都存在从
$\overline{\mathfrak C}_1$ 到 $\overline{\mathfrak C}_2$ 的
$A_K$ 同态 $\overline{\mathfrak F}^{(K)}$，它在
$\overline{\mathfrak f}$ 之上与 $\overline{\mathrm{Incl}}$、
$\overline{\mathrm{Eval}}$ 相容。此外，把
$\overline{\mathfrak F}^{(K)}$ 看成 $A_{K-1}$ 同态时，它与
$\overline{\mathfrak F}^{(K-1)}$ 相同。
\end{proposition}

\begin{proof}
对 $K$ 作归纳。$K=1$ 的情形已经由引理 4.4.2 处理。假设命题
4.4.11 对 $K-1$ 已经成立。于是
$\overline{\mathfrak F}^{(K-1)}$ 诱导余代数同态
\[
 \widehat{\overline{\mathfrak F}}^{(K-1)}\colon
 B\overline{\mathfrak C}_1[1]
 \longrightarrow B\overline{\mathfrak C}_2[1].
\]
考虑
\begin{equation}
 \widehat{\overline d}
 \circ\widehat{\overline{\mathfrak F}}^{(K-1)}
 -\widehat{\overline{\mathfrak F}}^{(K-1)}
 \circ\widehat{\overline d}
\tag{4.4.12}\label{eq:4-4-12}
\end{equation}
在 $B_{1\cdots K}\overline{\mathfrak C}_1[1]$ 上的限制，并把它记作
$\overline{\operatorname{Err}}_K$。

\begin{lemma}{4.4.13}\label{lem:4-4-13}
$\overline{\operatorname{Err}}_K$ 在
$B_{1\cdots K-1}\overline{\mathfrak C}_1[1]$ 上的限制为零；其像包含于
$B_1\overline{\mathfrak C}_2[1]=\overline{\mathfrak C}_2[1]$。
\end{lemma}

\begin{proof}
容易看出，$\widehat{\overline{\mathfrak F}}^{(K-1)}$ 在
$B_{1\cdots K-1}\overline{\mathfrak C}_1[1]$ 上的限制是
$\widehat{\overline{\mathfrak F}}_{1\cdots K-1}$。第一个断言由引理
4.4.10 得到。

$\widehat{\overline{\mathfrak F}}^{(K-1)}$ 在
\[
 \operatorname{Hom}\bigl(
 B_{2\cdots K}\overline{\mathfrak C}_1[1],
 B_{2\cdots K}\overline{\mathfrak C}_2[1]
 \bigr)
\]
上诱导 $\widehat{\overline{\mathfrak F}}_{2\cdots K}$。第二个断言也由
引理 4.4.10 得到。
\end{proof}

由引理 4.4.13，可以把
\[
 \overline{\operatorname{Err}}_K
 \in\operatorname{Hom}\bigl(
 B_K\overline{\mathfrak C}_1[1],
 \overline{\mathfrak C}_2[1]
 \bigr)
\]
看作一个同态。上边界算子 $\overline{\mathfrak m}_1$ 诱导
\[
 \delta_1\colon
 \operatorname{Hom}\bigl(
 B_K\overline{\mathfrak C}_1[1],
 \overline{\mathfrak C}_2[1]
 \bigr)
 \longrightarrow
 \operatorname{Hom}\bigl(
 B_K\overline{\mathfrak C}_1[1],
 \overline{\mathfrak C}_2[1]
 \bigr),
\]
其定义为
\[
 \delta_1(A)
 =\overline{\mathfrak m}_1\circ A
 -(-1)^{\deg A}A\circ\widehat{\overline{\mathfrak m}}_1.
\]

\begin{lemma}{4.4.14}\label{lem:4-4-14}
$\delta_1\overline{\operatorname{Err}}_K=0$。
\end{lemma}

\begin{proof}
注意，$\widehat{\overline d}$ 在
$\operatorname{Hom}(B\overline{\mathfrak C}_1[1],
B\overline{\mathfrak C}_2[1])$ 上诱导上边界算子
\[
 \phi\longmapsto
 \widehat{\overline d}(\phi)
 =\widehat{\overline d}\circ\phi
 -(-1)^{\deg\phi}\phi\circ\widehat{\overline d}.
\]
按定义，$\overline{\operatorname{Err}}_K$ 是
$\widehat{\overline d}
(\widehat{\overline{\mathfrak F}}^{(K-1)})$ 在
$B_K\overline{\mathfrak C}_1[1]$ 上的限制。因此，由引理 4.4.13，
$\widehat{\overline d}
(\widehat{\overline d}
(\widehat{\overline{\mathfrak F}}^{(K-1)}))=0$ 在
$B_K\overline{\mathfrak C}_1[1]$ 上的限制就是
$\delta_1\overline{\operatorname{Err}}_K$。引理得证。
\end{proof}

\begin{lemma}{4.4.15}\label{lem:4-4-15}
在
$\operatorname{Hom}(B_K\overline{\mathfrak C}_1[1],
\overline{\mathfrak C}_2[1])$ 中，有
$\overline{\operatorname{Err}}_K\in\operatorname{Im}\delta_1$。
\end{lemma}

\begin{proof}
令
\[
 \widehat{\overline{\mathfrak f}}^{(K-1)}\colon
 B\overline C_1[1]\longrightarrow B\overline C_2[1]
\]
为 $\widehat{\overline{\mathfrak f}}_1,\ldots,
\widehat{\overline{\mathfrak f}}_{K-1}$ 所诱导的余代数同态。
由于 $\overline{\mathfrak f}$ 是 $A_\infty$ 同态，$A_\infty$ 公式给出
\[
 \widehat{\overline d}
 \circ\widehat{\overline{\mathfrak f}}^{(K-1)}
 -\widehat{\overline{\mathfrak f}}^{(K-1)}
 \circ\widehat{\overline d}
 +\delta_1(\overline{\mathfrak f}_K)=0
\]
于 $B_{1\cdots K}\overline C_1[1]$ 上。这里
\[
 \delta_1\colon
 \operatorname{Hom}(B_K\overline C_1[1],\overline C_2[1])
 \longrightarrow
 \operatorname{Hom}(B_K\overline C_1[1],\overline C_2[1])
\]
由 $\overline{\mathfrak m}_1$ 诱导。

利用 $\overline{\mathfrak F}^{(K-1)}$ 与
$\overline{\mathrm{Incl}}$ 的相容性，在
$\operatorname{Hom}(B_K\overline C_1[1],
\overline{\mathfrak C}_2[1])$ 中计算得
\translatornote{原书在前一段把
$B_{1\cdots K}\overline C_1[1]$ 漏写为
$B_{1\cdots K}\overline C_1$，又在本句把固定的 $K$ 写成小写 $k$；
两处均由本小节的定义域及紧随的式 (4.4.16) 唯一确定，译文据此改正。}：
\begin{equation}
 \begin{aligned}
 \overline{\operatorname{Err}}_K
 \circ\widehat{\overline{\mathrm{Incl}}}
 &={}
 \Bigl(
 \widehat{\overline{\mathfrak m}}
 \circ\widehat{\overline{\mathfrak F}}^{(K-1)}
 -\widehat{\overline{\mathfrak F}}^{(K-1)}
 \circ\widehat{\overline d}
 \Bigr)
 \circ\widehat{\overline{\mathrm{Incl}}}\\
 &=\overline{\mathrm{Incl}}_1\circ
 \Bigl(
 \overline{\mathfrak m}
 \circ\widehat{\overline{\mathfrak f}}^{(K-1)}
 -\widehat{\overline{\mathfrak f}}^{(K-1)}
 \circ\widehat{\overline d}
 \Bigr)\\
 &=-\delta_1\bigl(
 \overline{\mathrm{Incl}}_1
 \circ\overline{\mathfrak f}_K
 \bigr).
 \end{aligned}
\tag{4.4.16}\label{eq:4-4-16}
\end{equation}
由于 $\overline{\mathrm{Incl}}_1$ 是上链同伦等价，引理随即成立。
\end{proof}

引理 4.4.15 蕴含存在
\[
 \overline{\mathfrak F}'_K
 \in\operatorname{Hom}\bigl(
 B_K\overline{\mathfrak C}_1[1],
 \overline{\mathfrak C}_2[1]
 \bigr),
 \qquad
 \delta_1(\overline{\mathfrak F}'_K)
 +\overline{\operatorname{Err}}_K=0.
\]

\begin{lemma}{4.4.17}\label{lem:4-4-17}
还可以选取 $\overline{\mathfrak F}'_K$，使得
\[
 \overline{\mathfrak F}'_K
 \circ\widehat{\overline{\mathrm{Incl}}}_1
 =\overline{\mathrm{Incl}}_1
 \circ\overline{\mathfrak f}_K.
\]
\end{lemma}

\begin{proof}
先任取满足
$\delta_1(\overline{\mathfrak F}''_K)
+\overline{\operatorname{Err}}_K=0$ 的
$\overline{\mathfrak F}''_K$。由 \eqref{eq:4-4-16}，
\[
 \delta_1\Bigl(
 \overline{\mathfrak F}''_K
 \circ\widehat{\overline{\mathrm{Incl}}}
 -\overline{\mathrm{Incl}}_1
 \circ\overline{\mathfrak f}_K
 \Bigr)=0.
\]
存在
\[
 \begin{aligned}
 \alpha&\in\operatorname{Hom}\bigl(
 B_K\overline{\mathfrak C}_1[1],
 \overline{\mathfrak C}_2[1]\bigr),\\
 \beta&\in\operatorname{Hom}\bigl(
 B_K\overline C_1[1],
 \overline{\mathfrak C}_2[1]\bigr),
 \end{aligned}
\]
使得 $\delta_1\alpha=0$，并且
\[
 \overline{\mathfrak F}''_K
 \circ\widehat{\overline{\mathrm{Incl}}}
 -\overline{\mathrm{Incl}}_1
 \circ\overline{\mathfrak f}_K
 =\alpha\circ\widehat{\overline{\mathrm{Incl}}}
 +\delta_1\beta.
\]
存在 $\widetilde\beta$，使得
$\widetilde\beta\circ\widehat{\overline{\mathrm{Incl}}}=\beta$。置
\[
 \overline{\mathfrak F}'_K
 =\overline{\mathfrak F}''_K-\alpha-\delta_1\widetilde\beta.
\]
容易检验 $\overline{\mathfrak F}'_K$ 具有所需性质。
\end{proof}

按定义，$\overline{\mathfrak F}^{(K-1)}$ 与
$\overline{\mathfrak F}'_K$ 共同定义了一个与
$\overline{\mathrm{Incl}}$ 相容的 $A_K$ 同态。

\begin{lemma}{4.4.18}\label{lem:4-4-18}
可以把 $\overline{\mathfrak F}'_K$ 修改为
$\overline{\mathfrak F}_K$，使得由
$\overline{\mathfrak F}^{(K-1)}$ 与 $\overline{\mathfrak F}_K$ 得到的
$A_K$ 同态 $\overline{\mathfrak F}^{(K)}$，在
$\overline{\mathfrak f}$ 之上与 $\overline{\mathrm{Incl}}$、
$\overline{\mathrm{Eval}}$ 相容。
\end{lemma}

\begin{proof}
先用 $\overline{\mathfrak F}'_K$ 得到 $A_K$ 同态
$\overline{\mathfrak F}^{(K)\prime}$。置
\begin{equation}
 \overline{\operatorname{Err}}_K^{(s_0)}
 =\overline{\mathrm{Eval}}_{s=s_0}
 \circ\overline{\mathfrak F}^{(K)\prime}
 -\overline{\mathfrak f}
 \circ\widehat{\overline{\mathrm{Eval}}}_{s=s_0},
 \qquad s_0=0,1.
\tag{4.4.19}\label{eq:4-4-19}
\end{equation}
由归纳假设，$\overline{\mathfrak F}^{(K-1)}$ 与
$\overline{\mathrm{Eval}}$ 相容，故
$\overline{\operatorname{Err}}_K^{(s_0)}$ 在
$B_{1\cdots K-1}\overline{\mathfrak C}_1[1]$ 上的限制为零。而且，它在
$B_{1\cdots K}\overline{\mathfrak C}_1[1]$ 上的限制之像包含于
$B_1\overline C_2[1]=\overline C_2[1]$。也就是说，
\[
 \bigl(
 \overline{\operatorname{Err}}_K^{(0)},
 \overline{\operatorname{Err}}_K^{(1)}
 \bigr)
 \in\operatorname{Hom}\bigl(
 B_K\overline{\mathfrak C}_1[1],
 \overline C_2[1]\oplus\overline C_2[1]
 \bigr).
\]
由于 $\overline{\mathfrak F}^{(K)\prime}$ 与
$\overline{\mathfrak f}$ 都是 $A_K$ 同态，
\[
 \delta_1\bigl(
 \overline{\operatorname{Err}}_K^{(0)},
 \overline{\operatorname{Err}}_K^{(1)}
 \bigr)=0.
\]
此外，由与 $\overline{\mathrm{Incl}}$ 的相容性，
\[
 \overline{\operatorname{Err}}_K^{(s_0)}
 \circ\widehat{\overline{\mathrm{Incl}}}=0.
\]
把引理 4.4.3 用于
$A=(\overline{\operatorname{Err}}_K^{(0)},
\overline{\operatorname{Err}}_K^{(1)})$，得到
\[
 \bigl(
 \overline{\operatorname{Cor1}}_K^{(0)},
 \overline{\operatorname{Cor1}}_K^{(1)}
 \bigr)
 \in\operatorname{Hom}\bigl(
 B_K\overline{\mathfrak C}_1[1],
 \overline C_2[1]\oplus\overline C_2[1]
 \bigr),
\]
使得
\begin{align}
 \delta_1\bigl(
 \overline{\operatorname{Cor1}}_K^{(0)},
 \overline{\operatorname{Cor1}}_K^{(1)}
 \bigr)
 &=\bigl(
 \overline{\operatorname{Err}}_K^{(0)},
 \overline{\operatorname{Err}}_K^{(1)}
 \bigr),
 \tag{4.4.20.1}\label{eq:4-4-20-1}\\
 \overline{\operatorname{Cor1}}_K^{(s_0)}
 \circ\widehat{\overline{\mathrm{Incl}}}&=0.
 \tag{4.4.20.2}\label{eq:4-4-20-2}
\end{align}
利用性质 (4.2.2.4)，可找到
\[
 \overline{\operatorname{Cor2}}_K
 \in\operatorname{Hom}\bigl(
 B_K\overline{\mathfrak C}_1[1],
 \overline{\mathfrak C}_2[1]
 \bigr),
\]
使得
\begin{align}
 (\overline{\mathrm{Eval}}_{s=s_0})_1
 \circ\overline{\operatorname{Cor2}}_K
 &=\overline{\operatorname{Cor1}}_K^{(s_0)},
 \tag{4.4.21.1}\label{eq:4-4-21-1}\\
 \overline{\operatorname{Cor2}}_K
 \circ\widehat{\overline{\mathrm{Incl}}}&=0.
 \tag{4.4.21.2}\label{eq:4-4-21-2}
\end{align}
置
\[
 \overline{\mathfrak F}_K
 =\overline{\mathfrak F}'_K-\overline{\operatorname{Cor2}}_K.
\]
于是 \eqref{eq:4-4-20-1}、\eqref{eq:4-4-21-1}、
\eqref{eq:4-4-21-2} 蕴含
$\overline{\mathfrak F}_K$ 具有所需性质。
\end{proof}

至此完成命题 4.4.11 的证明。
\end{proof}

命题 4.4.11 立即蕴含定理 4.2.34 的无滤过版本。
